\chapter{The Shuffle Framework and the Scheduler}
\label{ch:scheduler}

\begin{chapterintro}
\Cref{ch:shuffle} treated the shuffle as a wall. This chapter builds the wall:
the shuffle manager and its writers, the map output tracker that tells reduce
tasks where to fetch from, the DAG scheduler that cuts a plan into stages, the
task scheduler that places tasks on executors, and the block transfer service
that moves the bytes. It is the machinery behind every stage boundary and every
straggler.
\end{chapterintro}

\section{The shuffle manager and shuffle writers}
\label{sec:sched-shuffle-manager}

Every shuffle on a modern Spark version goes through \texttt{SortShuffleManager}
--- the separate hash-based shuffle manager from early Spark versions was
removed once sort-based shuffle's one-file-plus-index-per-task layout
(\Cref{sec:shuffle-write}) proved strictly better at scale. What varies is
which \emph{writer} a given shuffle uses underneath that one manager, chosen
automatically from the shuffle's own properties:

\begin{description}
  \item[\texttt{BypassMergeSortShuffleWriter}] fires when the shuffle has no
        map-side aggregation and the partition count is below
        \conf{spark.shuffle.sort.bypassMergeThreshold} (200 by default) ---
        cheap enough, at that few partitions, to just write one file per
        partition directly and concatenate them, skipping the sort step
        entirely.
  \item[\texttt{UnsafeShuffleWriter}] (the "serialized" path) fires when the
        serializer supports relocating serialized values without
        deserializing them (Kryo does, under most configurations) and there
        is no map-side aggregation --- it sorts the same compact
        \emph{(partition-id, pointer)} arrays \Cref{sec:tungsten-cache-aware}
        described, spilling and merging already-serialized bytes with no
        object ever reconstructed mid-sort.
  \item[\texttt{SortShuffleWriter}] is the general fallback --- a real
        record-by-record sort by partition (and key, if the shuffle
        aggregates) --- used whenever the other two writers' preconditions
        do not hold, at correspondingly higher cost.
\end{description}

The upshot: a shuffle you did not particularly optimize is not necessarily
running the expensive general path --- Spark already picks the cheapest
writer whose preconditions hold for that specific shuffle.

\section{The map output tracker}
\label{sec:sched-map-output-tracker}

A reduce task needs to know which executors hold which map task's output
before it can fetch anything --- that registry is the \emph{map output
tracker}. The driver's copy, \texttt{MapOutputTrackerMaster}, is the single
authoritative record, updated every time a map task finishes (with its
output location) and every time an executor is lost (with that executor's
outputs invalidated, which is precisely the trigger behind
\Cref{sec:failure-signatures}'s stage-resubmission behavior). Executors hold
a read-only, periodically-refreshed cache of the same information,
\texttt{MapOutputTrackerWorker}, so that a reduce task can resolve fetch
locations locally most of the time without an RPC to the driver for every
single task.

\section{The DAG scheduler and the task scheduler}
\label{sec:sched-dag-task}

Two schedulers do two different jobs, one layered on top of the other. The
\textbf{DAG scheduler} operates on RDD lineage: it walks the dependency
graph backward from an action, cutting a new stage boundary at every wide
dependency (\Cref{sec:shape-deps}), and submits stages to run only
once their upstream stages have finished --- this is the component that
turns one query into the job/stage structure \Cref{ch:shape} first
introduced, and the one that decides, on a \texttt{FetchFailedException},
exactly which upstream stage needs to re-run. The \textbf{task scheduler}
(\texttt{TaskSchedulerImpl}) operates one level down, inside an already-submitted
stage: given a set of runnable tasks and the current executors, it decides
\emph{which} task runs on \emph{which} executor, honoring the locality
preferences below and retrying a task, up to
\conf{spark.task.maxFailures}, on the executor it failed on before giving
up on that executor for that task.

\section{RPC and block transfer}
\label{sec:sched-rpc-transfer}

Coordination messages (task assignment, status updates, the map output
tracker's queries) travel over Spark's internal RPC layer, built on Netty;
the actual shuffle bytes travel over a separate path, the
\texttt{NettyBlockTransferService}, which is what a reduce task's fetch
request and an executor's block-serving response both ultimately go
through --- the same service the external shuffle service
(\Cref{sec:shuffle-external}) reimplements as a standalone process, so that
block-serving keeps working with no executor JVM behind it at all.

\section{Locality and speculation}
\label{sec:sched-locality-speculation}

Placing a task on the executor that already holds its input data avoids a
network read entirely; the task scheduler ranks candidate executors by
locality level --- \texttt{PROCESS\_LOCAL} (same JVM), \texttt{NODE\_LOCAL}
(same node, different executor), \texttt{RACK\_LOCAL}, then \texttt{ANY} ---
and will wait briefly (\conf{spark.locality.wait}, 3s by default) for a
better-locality slot to free up before falling back to a worse one. On this
stack's typical setup (compute and storage separated, GCS rather than
co-located HDFS), locality beyond \texttt{ANY} rarely applies to the initial
read regardless of this waiting, since there is no node the data is
particularly local to; it still matters for shuffle reads, where a task
scheduled on the same node as (some of) its shuffle input avoids that
portion of the network transfer.

Speculation (\Cref{sec:failure-signatures}) is the task scheduler's own
mechanism for the case locality cannot fix: once a task has run past
\conf{spark.speculation.multiplier} times the median duration for its stage,
the scheduler launches a second, speculative copy on a different executor
and accepts whichever attempt's result lands first, discarding the other ---
a scheduling-level decision independent of the DAG scheduler's stage-level
view, since a straggler is a single task problem, not a stage problem.

\section{Summary}
\label{sec:sched-summary}

One shuffle manager, three writers chosen by preconditions rather than
configuration you set directly: bypass for few partitions, serialized-sort
when the serializer allows it, plain sort otherwise. The map output tracker
is the registry that makes fetching possible and that a lost executor
invalidates, triggering the DAG scheduler's stage resubmission. The DAG
scheduler owns stage-level lineage and boundaries; the task scheduler owns
task placement within a stage, locality, and per-task retry, with
speculation as its answer to a straggler that locality alone cannot explain.
Both RPC coordination and the shuffle bytes themselves travel over the same
underlying Netty-based transport, just on logically separate paths.

\exercisehead
\begin{exercises}
\item Trace one wide query from \texttt{collect()} to task completion, naming
      each scheduler component it passes through.
\item Given a shuffle with 50 partitions and no aggregation, predict which
      shuffle writer handles it, then confirm from
      \conf{spark.shuffle.sort.bypassMergeThreshold}.
\item Kill an executor mid-stage and, from the event log, find the map
      output tracker invalidation and the DAG scheduler's resubmitted
      stage attempt.
\end{exercises}
